\chapter{Formal Metric Definitions}\label{app:metrics}

This appendix gives closed-form definitions for the taxonomy-quality
and structural-characterisation metrics named in
Section~\ref{sec:exp-metrics} and used in
Chapter~\ref{ch:results}. It does not cover the arena/ranking metrics
(Bradley--Terry scores, Kendall $\tau$, Spearman $\rho$,
pairwise-winner accuracy, top-$k$ Jaccard), which are already
formalised in Section~\ref{sec:meth-bt} and
Section~\ref{sec:meth-validation}.

\section{Spherical Silhouette Coefficient}\label{sec:app-metrics-silhouette}

Standard silhouette coefficients use Euclidean distance; for
L2-normalised embeddings on the unit hypersphere, cosine-based
angular distance is the geometrically appropriate analogue. For two
unit vectors $x, y \in \mathcal{S}^{d-1}$, the angular distance is
\[
d_\theta(x, y) = \arccos(x^\top y).
\]
For each query embedding $x$ with assigned node $A$ and nearest
neighbouring node $B \neq A$ (by centroid distance),
\[
a(x) = \arccos(x^\top \mu_A), \qquad
b(x) = \min_{B \neq A} \arccos(x^\top \mu_B), \qquad
s(x) = \frac{b(x) - a(x)}{\max(a(x), b(x))},
\]
with $s(x) = 0$ when $\max(a(x), b(x)) = 0$. The reported Spherical
Silhouette is the mean $\bar{S} = \frac{1}{n}\sum_i s(x_i)$, with
range $[-1, 1]$; values close to $1$ indicate queries tightly aligned
with their own node centroid and well separated from the nearest
neighbouring node.

\section{Total Dasgupta Cost}\label{sec:app-metrics-dasgupta}

The full-tree Dasgupta cost complements the per-split
chance-corrected separation score used during construction
(Appendix~\ref{app:dag-formalism}, \S\ref{sec:meth-split}). For similarity
matrix $W$ with $w_{ij} = \max(x_i^\top x_j, 0)$ (cosine similarity,
clipped to remove negative weights arising from MRL prefix
mismatches) and completed hierarchy $T$,
\[
\text{cost}(T) = \sum_{i \neq j} w_{ij} \, |\text{Leaves}(\text{LCA}_T(i,j))|,
\]
where $\text{LCA}_T(i,j)$ is the lowest common ancestor of queries
$i$ and $j$ in $T$ and $|\text{Leaves}(\cdot)|$ counts the leaves
descending from it \parencite{dasgupta2016cost}. Similar query pairs
that split deep in the hierarchy incur a low penalty; pairs that
split at the root are penalised by the full leaf count. A lower Total
Dasgupta Cost indicates a more structurally coherent hierarchy; unlike
the per-split chance-corrected separation acceptance test of
Appendix~\ref{app:dag-formalism} (\S\ref{sec:meth-split})---which is
not a Dasgupta quantity---it is a genuine
whole-tree Dasgupta-cost measure, computed on a fixed dataset and
comparable across configurations and
baselines (Section~\ref{sec:exp-baselines}).

\section{Normalised Sackin Index}\label{sec:app-metrics-sackin}

The Sackin index $S(T) = \sum_{\ell \in \text{Leaves}} \text{depth}(\ell)$
sums leaf depths; it is a standard tree-balance statistic in
phylogenetics. To remove its dependence on leaf count, the
\textbf{Normalised Sackin Index} is the mean leaf depth,
\[
\tilde{S}(T) = \frac{1}{|\text{Leaves}|} \sum_{\ell \in \text{Leaves}} \text{depth}(\ell).
\]
A low $\tilde{S}(T)$ indicates a flat hierarchy that splits early and
close to the root; a high $\tilde{S}(T)$ indicates a deep, uneven
hierarchy in which some domains decompose much further than others.

\section{Average Match Count}\label{sec:app-metrics-avgmatch}

In the polyhierarchical DAG (Section~\ref{sec:arch-dag}), a query can
be assigned to more than one leaf. Writing $\mathcal{L}(q)$ for the
set of leaves query $q$ is assigned to,
\[
\text{AvgMatch} = \frac{1}{N} \sum_{i=1}^{N} |\mathcal{L}(q_i)|.
\]
$\text{AvgMatch} = 1.0$ means the taxonomy behaves as a strict tree
for routing purposes; $\text{AvgMatch} > 1.0$ confirms that
soft-routing and bridge insertion (Appendix~\ref{app:dag-formalism},
\S\ref{sec:meth-topology})
are actively producing multi-membership.

\section{Weighted Leaf Purity}\label{sec:app-metrics-wlp}

Weighted Leaf Purity (WLP) is not a standard named metric in the
taxonomy-induction literature; it is defined here as a size-weighted
average of per-leaf label purity against the MMLU-Pro ground-truth
label. For leaf $\ell$ with member set $Q_\ell$ and ground-truth label
function $\text{gt}(\cdot)$,
\[
\text{purity}(\ell) = \max_{c} \frac{|\{q \in Q_\ell : \text{gt}(q) = c\}|}{|Q_\ell|},
\qquad
\text{WLP} = \sum_{\ell \in \text{Leaves}} \frac{|Q_\ell|}{N} \, \text{purity}(\ell).
\]
Because WLP flat-tests each leaf directly rather than through a
DAG lowest-common-ancestor traversal, it is a useful complement to
the LCA-sensitive Dendrogram Purity of
Section~\ref{sec:app-metrics-dendrogram-purity}.

\section{DAG Dendrogram Purity}\label{sec:app-metrics-dendrogram-purity}

Standard dendrogram purity is defined over a tree, where the lowest
common ancestor (LCA) of two leaves is unique. In a DAG, a node can
have more than one parent (Appendix~\ref{app:dag-formalism},
\S\ref{sec:meth-topology}), so the
LCA of two queries is not in general unique. The implementation
resolves this with the \emph{shallowest} LCA taken over the full
DAG (a conservative bound), but tests subtree purity only over the
tree-skeleton (primary-edge) structure---cross-links are followed
when locating the common ancestor but not when enumerating the
purity-test subtree's members. This hybrid avoids the metric
collapsing toward its floor value under heavy cross-linking. Because
the shallowest-LCA convention means same-label queries split across
distant subtrees still resolve to a high, impure ancestor, Dendrogram
Purity should be read alongside the flat Weighted Leaf Purity of
Section~\ref{sec:app-metrics-wlp}, which is insensitive to LCA
placement.

\section{Routing Expected Calibration Error}\label{sec:app-metrics-ece}

Routing ECE tests whether the soft-routing distribution's confidence
is calibrated against its accuracy, evaluated at \emph{domain}
granularity. For each query, its normalized leaf membership weights
(the real trickle weights of Appendix~\ref{app:dag-formalism},
\S\ref{sec:meth-routing}) are aggregated onto each leaf's depth-1
domain ancestors, yielding a per-query distribution over domain
categories in the same label space as the MMLU-Pro category strings;
the query's \emph{confidence} is the maximum domain mass and its
prediction is correct when that argmax domain matches the query's
MMLU-Pro category. Queries are grouped
into $B$ equal-width confidence bins; for bin $b$ with queries
$\mathcal{B}_b$, $\text{acc}(\mathcal{B}_b)$ is the fraction of
correct predictions in the bin and $\text{conf}(\mathcal{B}_b)$ is
the bin's mean confidence. Then
\[
\text{ECE} = \sum_{b=1}^{B} \frac{|\mathcal{B}_b|}{N} \left| \text{acc}(\mathcal{B}_b) - \text{conf}(\mathcal{B}_b) \right|,
\]
the standard expected-calibration-error construction adapted from
temperature-scaling calibration analysis to the domain-aggregated
trickle-routing distribution. Lower Routing ECE
indicates a routing distribution whose confidence is a trustworthy
proxy for correctness.
